\documentclass[11pt]{article}
\usepackage[margin=1in]{geometry}
\usepackage[hidelinks]{hyperref}
\usepackage{longtable}
\usepackage{array}
\usepackage{enumitem}

\title{Build-a-Bot Stadium\\As-Built Summary}
\author{Project Snapshot}
\date{18 March 2026}

\begin{document}
\maketitle

\section{Purpose and Scope}
This document summarizes the current \emph{as-built} state of Build-a-Bot Stadium at a practical, implementation-aligned level. It is intended as a concise technical orientation for contributors and reviewers who need to understand what the platform currently does, how it is organized, and where the main gaps remain.

The summary focuses on implemented capabilities and operational characteristics rather than design aspirations.

\section{System Overview}
Build-a-Bot Stadium is currently implemented as a single Go runtime with two primary external surfaces:
\begin{itemize}[leftmargin=*]
  \item A TCP bot protocol endpoint (default \texttt{:8080})
  \item An HTTP dashboard and viewer endpoint (default \texttt{:3000})
\end{itemize}

The platform now supports three workload styles:
\begin{itemize}[leftmargin=*]
  \item Built-in competitive games (for example Connect4)
  \item Built-in single-agent episodic environments (for example Gridworld)
  \item Optional process-based external game plugins discovered at runtime
\end{itemize}

Primary references:
\begin{itemize}[leftmargin=*]
  \item \texttt{README.md}
  \item \texttt{ARCHITECTURE.md}
\end{itemize}

\section{Implemented Feature Set}
\subsection{Arena Runtime and State Management}
The in-memory manager in \texttt{stadium/} acts as the source of truth for sessions, arenas, bot profiles, and archived match records. Arena lifecycle, participant attachment, outcomes, and dashboard snapshot broadcasting are handled centrally.

Notable implemented behavior:
\begin{itemize}[leftmargin=*]
  \item One-player and two-player arena activation support via game policy interfaces
  \item Move clock and handicap policy handling based on game capabilities
  \item Match archival with replay-oriented metadata
  \item Watchdog-based arena cleanup
\end{itemize}

Primary references:
\begin{itemize}[leftmargin=*]
  \item \texttt{ARCHITECTURE.md}
  \item \texttt{stadium/manager.go} and related files in \texttt{stadium/}
\end{itemize}

\subsection{Bot Command Protocol (TCP)}
A newline-delimited command protocol is exposed for bot clients. Core commands include registration, arena creation/joining, move submission, observation, and disconnect.

Current protocol supports dynamic game arguments in \texttt{CREATE}, enabling consistent invocation for built-in and plugin games.

Primary references:
\begin{itemize}[leftmargin=*]
  \item \texttt{PROTOCOL.md}
  \item \texttt{cmd/bbs-server/main.go}
\end{itemize}

\subsection{Dashboard and Viewer Surfaces}
The dashboard provides operational visibility and control over runtime state, including owner-scoped and admin-scoped control paths.

Implemented dashboard capabilities include:
\begin{itemize}[leftmargin=*]
  \item Live state streaming via SSE
  \item Owner token flow (register bot, create/join/leave/eject)
  \item Admin control flow protected by \texttt{BBS\_DASHBOARD\_ADMIN\_KEY}
  \item Dynamic arena creation forms driven by live game catalog metadata
  \item Live and replay viewer pages
\end{itemize}

Recent plugin-oriented UX additions include a plugin discovery diagnostics panel showing enabled state, scan directory, and per-manifest load/skip status.

Primary references:
\begin{itemize}[leftmargin=*]
  \item \texttt{cmd/bbs-server/dashboard.go}
  \item \texttt{cmd/bbs-server/templates/index.html}
  \item \texttt{cmd/bbs-server/templates/dashboard-state.html}
  \item \texttt{DASHBOARD\_REMOTE\_QUICKSTART.md}
\end{itemize}

\subsection{Process-Based Game Plugin System}
A process plugin model is now scaffolded and operational behind feature flags:
\begin{itemize}[leftmargin=*]
  \item \texttt{BBS\_ENABLE\_GAME\_PLUGINS=true}
  \item \texttt{BBS\_GAME\_PLUGIN\_DIR=/path/to/plugins}
\end{itemize}

Plugin manifests are discovered from \texttt{*.json} files and merged with built-in registrations into a single runtime catalog. Plugin games communicate with the host over JSONL RPC via stdin/stdout.

Implemented plugin author support:
\begin{itemize}[leftmargin=*]
  \item Shared protocol and helper runtime (\texttt{games/pluginapi})
  \item Reference plugin command (\texttt{cmd/bbs-game-counter-plugin})
  \item Dashboard-integrated arg schema rendering
  \item Discovery diagnostics in the dashboard
\end{itemize}

Primary references:
\begin{itemize}[leftmargin=*]
  \item \texttt{ARCHITECTURE.md}
  \item \texttt{PLUGIN\_AUTHORING.md}
  \item \texttt{games/pluginapi/protocol.go}
  \item \texttt{games/pluginapi/server.go}
  \item \texttt{games/plugin\_host.go}
  \item \texttt{games/plugin\_process\_game.go}
\end{itemize}

\subsection{Local Bridge Agent for Bot Authors}
\texttt{bbs-agent} provides a local JSONL bridge abstraction for bot logic, separating local policy code from direct BBS TCP protocol handling.

Implemented bridge model:
\begin{itemize}[leftmargin=*]
  \item Local socket endpoint and message contract v0.2
  \item \texttt{hello} / \texttt{welcome} / \texttt{turn} / \texttt{action} message loop
  \item Support for both competitive and environment-style usage patterns
\end{itemize}

Primary references:
\begin{itemize}[leftmargin=*]
  \item \texttt{BBS\_AGENT\_CONTRACT.md}
  \item \texttt{cmd/bbs-agent/README.md}
  \item \texttt{cmd/bbs-agent/main.go}
\end{itemize}

\subsection{Packaging, Deployment, and Release Automation}
Deployment guidance exists for both Docker and non-Docker TrueNAS workflows, and release automation publishes Linux binaries for amd64 and arm64.

Primary references:
\begin{itemize}[leftmargin=*]
  \item \texttt{deploy/truenas/README.md}
  \item \texttt{deploy/truenas-no-docker/README.md}
  \item \texttt{.github/workflows/release-bbs-server.yml}
\end{itemize}

\subsection{Validation and CI Status}
A dedicated manifest-validation command and workflow are present:
\begin{itemize}[leftmargin=*]
  \item \texttt{cmd/bbs-plugin-manifest-lint/main.go}
  \item \texttt{.github/workflows/validate-plugin-manifests.yml}
\end{itemize}

This improves plugin packaging hygiene by validating schema shape, field constraints, duplicates, and protocol compatibility.

\section{Current Gaps and Shortcomings}
The following are the most immediate gaps visible in the current implementation snapshot.

\subsection{Durability and Data Lifecycle}
\begin{itemize}[leftmargin=*]
  \item Runtime state is in-memory only (sessions, arenas, profiles, and match history).
  \item Process restart clears active runtime context and history.
\end{itemize}

Implication: suitable for local/home-lab operation and experimentation, but not yet a durable production service.

\subsection{Security and Trust Model}
\begin{itemize}[leftmargin=*]
  \item Bot channel remains plain TCP without built-in transport security.
  \item Identity and owner-token controls are bearer-style.
  \item Plugin trust assumes local trusted filesystem and trusted binaries.
\end{itemize}

Implication: deployment in untrusted networks requires external controls (for example TLS termination, network segmentation, host hardening).

\subsection{Protocol Consistency}
\begin{itemize}[leftmargin=*]
  \item Protocol documentation notes that some legacy command paths still emit plain text rather than normalized JSON envelopes.
\end{itemize}

Implication: client implementations need defensive parsing in edge paths until response shape is fully unified.

\subsection{Scalability and Concurrency Model}
\begin{itemize}[leftmargin=*]
  \item Manager state synchronization is intentionally coarse-grained (central mutex over key maps/counters).
\end{itemize}

Implication: simple and robust at current scale, but likely to become a throughput bottleneck under heavier concurrent load.

\subsection{Plugin Governance and Isolation}
\begin{itemize}[leftmargin=*]
  \item Plugin discovery and shape validation exist, but stronger controls are not yet present (for example plugin signing, attestation, quotas, sandbox policies).
  \item Health enforcement and richer lifecycle supervision are limited.
\end{itemize}

Implication: plugin extensibility is practical now, but governance and safety controls are early-stage.

\subsection{Automated Test Coverage}
Repository snapshot indicates no \texttt{*\_test.go} unit/integration test files currently present.

Implication: current confidence is build- and smoke-test-oriented, with limited regression guardrails from automated test suites.

\section{Suggested Near-Term Focus}
A pragmatic next phase that aligns with the current architecture would be:
\begin{enumerate}[leftmargin=*]
  \item Persist core runtime records (profiles, matches, key arena events) to durable storage.
  \item Normalize protocol outputs so all command responses consistently use one JSON envelope model.
  \item Introduce baseline automated tests for manager lifecycle, protocol handlers, and plugin host behavior.
  \item Add security hardening layers for production-like deployments (transport protection, token handling improvements, plugin controls).
  \item Expand plugin operational controls (health checks, richer diagnostics, controlled resource limits).
\end{enumerate}

\section{Document Map (Primary Reading Order)}
\begin{longtable}{|p{0.34\textwidth}|p{0.60\textwidth}|}
\hline
\textbf{Document} & \textbf{Primary Use} \\
\hline
\path{README.md} & High-level platform orientation, run instructions, plugin quickstart, and usage examples. \\
\hline
\path{ARCHITECTURE.md} & Runtime decomposition, component responsibilities, plugin model, and known boundaries. \\
\hline
\path{PROTOCOL.md} & TCP command protocol and response conventions for bot clients. \\
\hline
\path{BBS_AGENT_CONTRACT.md} & Local bridge JSONL contract for \texttt{bbs-agent} integrations. \\
\hline
\path{PLUGIN_AUTHORING.md} & Plugin manifest schema, local workflow, CI checks, and release checklist. \\
\hline
\path{DASHBOARD_REMOTE_QUICKSTART.md} & Remote operator workflow for connecting bots through dashboard-issued owner tokens. \\
\hline
\path{deploy/truenas/README.md} & Docker-based TrueNAS deployment and update operations. \\
\hline
\path{deploy/truenas-no-docker/README.md} & Non-Docker TrueNAS deployment and update operations. \\
\hline
\path{.github/workflows/release-bbs-server.yml} & Binary release automation (linux amd64 and arm64 artifacts). \\
\hline
\path{.github/workflows/validate-plugin-manifests.yml} & CI manifest validation workflow for plugin packaging quality. \\
\hline
\end{longtable}

\section{As-Built Conclusion}
Build-a-Bot Stadium has progressed from a narrow game host into a broader, extensible bot-arena platform with:
\begin{itemize}[leftmargin=*]
  \item multi-mode arena semantics,
  \item a functional process-plugin architecture,
  \item integrated dashboard-driven operations, and
  \item practical deployment and release pathways.
\end{itemize}

The core extension and operator workflows are now demonstrably in place. The most significant remaining work is concentrated in durability, protocol cleanup, security hardening, and automated regression coverage.

\end{document}
